% Generated draft: requires mathematical and source review.
\documentclass[11pt]{article}
\usepackage[margin=25mm]{geometry}
\usepackage{fontspec}
\setmainfont{Latin Modern Roman}
\usepackage{amsmath,amssymb,mathtools,unicode-math}
\setmathfont{Latin Modern Math}
\usepackage{enumitem,needspace,xurl,hyperref}
\setlength{\parindent}{0pt}
\setlength{\parskip}{5pt}
\setlength{\emergencystretch}{3em}
\newcommand{\N}{\mathbb{N}}
\newcommand{\Z}{\mathbb{Z}}
\newcommand{\Q}{\mathbb{Q}}
\newcommand{\R}{\mathbb{R}}
\newcommand{\C}{\mathbb{C}}
\newcommand{\F}{\mathbb{F}}
\newcommand{\A}{\mathbb{A}}
\newcommand{\PP}{\mathbb{P}}
\newcommand{\E}{\mathbb{E}}
\newcommand{\eps}{\varepsilon}
\newcommand{\dd}{\,\mathrm{d}}
\begin{document}
\section*{UnsolvedMath ID: 30003848}
\textbf{Source problem code:} catalogue entry
\subsection{Definitions and mathematical statement}
Let \(D=\Z[1/2]\). Thompson's group \(F\) is the group of increasing piecewise linear homeomorphisms of \([0,1]\) with finitely many breakpoints, all in \(D\), and slopes in \(\{2^k:k\in\Z\}\). Thompson's group \(T\) is the group of orientation-preserving piecewise linear homeomorphisms of \(\R/\Z\) with finitely many breakpoints in \(D/\Z\), slopes in \(\{2^k:k\in\Z\}\), and preserving \(D/\Z\). The notation \(F\times F\) denotes the direct product.

The supplied source asks whether \(F\) has any ``interesting proper nonabelian quasi-retracts,'' and specifically asks:
\begin{enumerate}
\item Is \(F\) a quasi-retract of \(T\)?
\item Is \(F\times F\) a quasi-retract of \(F\)?
\end{enumerate}
The supplied record does not define ``quasi-retract,'' specify the relevant embeddings, or give a mathematical criterion for ``interesting.'' Consequently, the following metric formulation is conditional on a convention that requires confirmation against the report.

For a finitely generated group \(G\) with word metric \(d_G\) and a finitely generated subgroup \(H\leq G\) with its own word metric \(d_H\), a coarse Lipschitz retraction means a map \(r:G\to H\) and constants \(L\geq1\), \(A,B\geq0\) such that
\[
 d_H(r(x),r(y))\leq Ld_G(x,y)+A
 \quad\text{for all }x,y\in G,
 \qquad
 d_H(r(h),h)\leq B
 \quad\text{for all }h\in H.
\]
Under the convention that existence of such a map makes \(H\) a quasi-retract of \(G\), the two questions ask for such retractions for the respective pairs. Natural candidate embeddings are \(F\hookrightarrow T\), obtained by identifying the endpoints of \([0,1]\), and \(F\times F\hookrightarrow F\), obtained by letting the two factors act on \([0,1/2]\) and \([1/2,1]\) after affine rescaling. Neither this convention nor these particular embeddings are specified in the supplied record; they are proposed interpretations, not confirmed source hypotheses. The broader question about ``interesting'' proper nonabelian quasi-retracts remains an informal request without a specified class of subgroups or a criterion for interest.
\subsection{Short English statement}
The source asks whether Thompson's group \(F\) has interesting proper nonabelian quasi-retracts, and separately whether \(F\) is a quasi-retract of \(T\) and whether \(F\times F\) is a quasi-retract of \(F\). A precise metric interpretation is supplied conditionally because the record does not define quasi-retract or specify embeddings. The supplied open status requires independent verification.
\subsection{Sources}
\begin{itemize}
\item Supplied local UnsolvedMath registry record, permanent ID \(30003848\), source problem code \texttt{catalogue entry}, in the Oberwolfach Reports Open Problems collection. No individual proposer is supplied; the proposed year is recorded as 2018.
\item Supplied source citation: \emph{Cohomological and Metric Properties of Groups of Homeomorphisms of R} (2019), Oberwolfach Reports, pp.\ 1579--1633, DOI \texttt{10.4171/owr/2018/26}. \url{https://doi.org/10.4171/owr/2018/26}.
\item Supplied literature source: \emph{Oberwolfach Report 26/2018}. \url{https://ems.press/content/serial-article-files/46748?nt=1}.
\item Literature assessment dated exactly \texttt{2026-08-21}: the supplied curation labels the problem open, states that the report poses both quasi-retract questions as open, and reports that no later primary resolution of either question was found. This is supplied dated curation, not an independently verified literature assessment; neither linked source was visited for this draft.
\item The registry labels the statement \texttt{corrected\_verified} and reports a manual semantic, typographic, and LaTeX review. These are supplied editorial assertions, not independent mathematical verification of this draft.
\end{itemize}
\end{document}
